\section{Discussion}
\label{sec:discussion}

\subsection{Comparison of the three approaches}
\Cref{tab:approaches} summarizes the strengths and weaknesses observed for the three families.

\begin{table}[H]
  \centering
  \small
  \begin{tabularx}{\linewidth}{l>{\raggedright\arraybackslash}X>{\raggedright\arraybackslash}X}
    \toprule
    \textbf{Approach} & \textbf{Strengths} & \textbf{Limitations} \\
    \midrule
    Traditional pipeline &
    Best accuracy (gradient boosting); fast to train (seconds); every preprocessing decision is explicit and testable; mature tooling for tuning and thresholding &
    Linear and distance-based models depend heavily on hand-made preprocessing; ensembles are harder to explain than a single model \\
    Deep learning &
    Learns representations of categorical variables end to end; beats every linear model without manual interaction terms; the learned embeddings are reusable features &
    Below the tree ensembles on this medium-sized table; many training hyper-parameters; slower (TabNet especially); less transparent \\
    Bayesian network &
    Explicit, inspectable assumptions; exact answers to any conditional query, in both predictive and diagnostic directions; native handling of missing features &
    Requires discretization, which loses information; CPT size grows exponentially with the number of parents; weakest classifier of the three families \\
    \bottomrule
  \end{tabularx}
  \caption{Strengths and limitations of the three model families on this problem.}
  \label{tab:approaches}
\end{table}

\subsection{What mattered most}
Three decisions had a large effect, and most others had almost none.
\begin{enumerate}
  \item \textbf{The model family.} Moving from a linear model to a tree ensemble raised the test PR-AUC from 0.75 to 0.90, more than every preprocessing choice combined.
  \item \textbf{Scaling, for distance- and margin-based models.} It decides whether k-NN and the SVM are usable at all.
  \item \textbf{The decision threshold.} For a fixed model, moving the threshold trades recall against precision far more than any re-weighting or resampling strategy does. Tuning it for F2 is a direct, transparent way to encode the higher cost of a missed defaulter.
\end{enumerate}
In contrast, the imputation method, the treatment of the seven impossible values, the scaler variant and the resampling method changed PR-AUC by less than the fold-to-fold noise.
Reporting these negative results matters: they show which parts of the pipeline deserve attention in practice.

\subsection{Choice of metric}
Accuracy would have been misleading for this problem.
For example, the random forest has a higher test accuracy than gradient boosting at their tuned thresholds, but a lower recall; a lender who cares about missed defaulters should prefer the latter.
PR-AUC ranks models by how well they separate defaulters across all thresholds, and F2 makes the cost asymmetry explicit when choosing one threshold.
The weight $\beta = 2$ is itself a business assumption: a lender with a different ratio between the cost of a default and the lost margin of a rejected applicant should choose a different $\beta$, or minimize the expected cost directly.

\subsection{Limitations and threats to validity}
\begin{itemize}
  \item \textbf{One dataset and one split.} All conclusions come from one public dataset with a single stratified test split. Cross-validation on the training set agrees with the test results, but the dataset's provenance (country, period, lender) is undocumented, so the absolute numbers may not transfer to other portfolios.
  \item \textbf{No temporal information.} Credit models are normally validated on later loans than they are trained on. The dataset has no dates, so a random split may slightly overstate performance under drift.
  \item \textbf{Limited search spaces.} The grids are deliberately small to keep \emph{Run all} within the limits of a free Colab session. The headroom experiment (\cref{tab:headroom}) suggests this does not change the conclusions for the best family, but other families might gain somewhat from wider searches.
  \item \textbf{Fixed choices.} The IQR factor $k = 1.5$, the four engineered features, the MLP training schedule and the Bayesian-network parent limit were set by convention rather than tuned.
  \item \textbf{Fairness.} The features include age, and home ownership acts as a proxy for wealth. A deployed credit model must be audited for disparate impact, which is outside the scope of this assignment.
  \item \textbf{Calibration.} Class weights and threshold tuning shift the predicted probabilities, so the scores should be calibrated before they are read as default probabilities, for example to price a loan.
\end{itemize}
